\documentclass[10pt,a4paper]{amsart}
\usepackage{amsmath,amssymb,amsthm,graphicx}
% Retained preamble for the Stacks Project chapter divisors.tex (Divisors).
% The chapter file pinned-original-source.tex carries only \input{preamble} on its first line, so no
% preamble text lives inside the chapter itself.  The Stacks Project uses one single shared preamble for
% every chapter; the snapshot stages it verbatim in sources/stacks-chow/preamble.tex.  The block below is
% copied mechanically, line for line, from that file's lines 14-18 (the Xy-pic setup) and lines 154-217
% (the native theorem declarations together with the mathematical macros).  Nothing was added, removed or
% redefined.  The wrapper class and the amsmath/amssymb/amsthm/graphicx base packages are supplied by the
% assembler, so the class/\externaldocument part of the shared preamble is not repeated here.
% For commutative diagrams we use Xy-pic
\usepackage[all]{xy}

% We use 2cell for 2-commutative diagrams.
\xyoption{2cell}
% Theorem environments.
%
\theoremstyle{plain}
\newtheorem{theorem}[subsection]{Theorem}
\newtheorem{proposition}[subsection]{Proposition}
\newtheorem{lemma}[subsection]{Lemma}

\theoremstyle{definition}
\newtheorem{definition}[subsection]{Definition}
\newtheorem{example}[subsection]{Example}
\newtheorem{exercise}[subsection]{Exercise}
\newtheorem{situation}[subsection]{Situation}

\theoremstyle{remark}
\newtheorem{remark}[subsection]{Remark}
\newtheorem{remarks}[subsection]{Remarks}

\numberwithin{equation}{subsection}

% Macros
%
\def\lim{\mathop{\mathrm{lim}}\nolimits}
\def\colim{\mathop{\mathrm{colim}}\nolimits}
\def\Spec{\mathop{\mathrm{Spec}}}
\def\Hom{\mathop{\mathrm{Hom}}\nolimits}
\def\Ext{\mathop{\mathrm{Ext}}\nolimits}
\def\SheafHom{\mathop{\mathcal{H}\!\mathit{om}}\nolimits}
\def\SheafExt{\mathop{\mathcal{E}\!\mathit{xt}}\nolimits}
\def\Sch{\mathit{Sch}}
\def\Mor{\mathop{\mathrm{Mor}}\nolimits}
\def\Ob{\mathop{\mathrm{Ob}}\nolimits}
\def\Sh{\mathop{\mathit{Sh}}\nolimits}
\def\NL{\mathop{N\!L}\nolimits}
\def\CH{\mathop{\mathrm{CH}}\nolimits}
\def\proetale{{pro\text{-}\acute{e}tale}}
\def\etale{{\acute{e}tale}}
\def\QCoh{\mathit{QCoh}}
\def\Ker{\mathop{\mathrm{Ker}}}
\def\Im{\mathop{\mathrm{Im}}}
\def\Coker{\mathop{\mathrm{Coker}}}
\def\Coim{\mathop{\mathrm{Coim}}}

% Boxtimes
%
\DeclareMathSymbol{\boxtimes}{\mathbin}{AMSa}{"02}

%
% Macros for moduli stacks/spaces
%
\def\QCohstack{\mathcal{QC}\!\mathit{oh}}
\def\Cohstack{\mathcal{C}\!\mathit{oh}}
\def\Spacesstack{\mathcal{S}\!\mathit{paces}}
\def\Quotfunctor{\mathrm{Quot}}
\def\Hilbfunctor{\mathrm{Hilb}}
\def\Curvesstack{\mathcal{C}\!\mathit{urves}}
\def\Polarizedstack{\mathcal{P}\!\mathit{olarized}}
\def\Complexesstack{\mathcal{C}\!\mathit{omplexes}}
% \Pic is the operator that assigns to X its picard group, usage \Pic(X)
% \Picardstack_{X/B} denotes the Picard stack of X over B
% \Picardfunctor_{X/B} denotes the Picard functor of X over B
\def\Pic{\mathop{\mathrm{Pic}}\nolimits}
\def\Picardstack{\mathcal{P}\!\mathit{ic}}
\def\Picardfunctor{\mathrm{Pic}}
\def\Deformationcategory{\mathcal{D}\!\mathit{ef}}

% --- appended 2026-10-12 (Chief): the project own definitions of the reference/slogan/history
% display environments, copied mechanically from the shared Stacks preamble lines 8-12
% (\usepackage{verbatim} plus the three \newenvironment lines, which the chapter preamble keeps
% verbatim from the project preamble). Nothing added, removed or redefined: these environments
% carry no mathematical content and the project itself typesets them as comments. Without these
% three lines a counted statement whose own text uses them fails with Environment slogan undefined.
\usepackage{verbatim}
\newenvironment{reference}{\comment}{\endcomment}
\newenvironment{slogan}{\comment}{\endcomment}
\newenvironment{history}{\comment}{\endcomment}
\usepackage{cleveref}
\crefname{originallemma}{original lemma}{original lemmas}
\Crefname{originallemma}{Original lemma}{Original lemmas}
\makeatletter
\newlabel{lemma-closed-subscheme-proj}{{0801}{source}{}{sourcepointer-1}{}}
\newlabel{lemma-closed-subscheme-proj@cref}{{[originallemma][0][]0801}{[1][0][]source}{}{}{}}
\makeatother
\crefname{originallemma}{original lemma}{original lemmas}
\Crefname{originallemma}{Original lemma}{Original lemmas}
\makeatletter
\newlabel{lemma-closed-subscheme-proj-finite}{{0802}{source}{}{sourcepointer-2}{}}
\newlabel{lemma-closed-subscheme-proj-finite@cref}{{[originallemma][0][]0802}{[1][0][]source}{}{}{}}
\makeatother
\crefname{originallemma}{original lemma}{original lemmas}
\Crefname{originallemma}{Original lemma}{Original lemmas}
\makeatletter
\newlabel{properties-lemma-quasi-coherent-colimit-finite-type}{{01PG}{source}{}{sourcepointer-3}{}}
\newlabel{properties-lemma-quasi-coherent-colimit-finite-type@cref}{{[originallemma][0][]01PG}{[1][0][]source}{}{}{}}
\makeatother
\title{A closed subscheme of a relative Proj away from a quasi-compact open is cut out by a finite type graded piece}
\author{The Stacks Project Authors}
\date{stacks-project:divisors.tex:local-snapshot-20261012; GNU FDL 1.2 or later}
\begin{document}
\maketitle
\noindent\textit{Editorial note.} All delivered mathematical text is the author's own, copied from the pinned Stacks Project chapter file pinned-original-source.tex; no statement or proof step was rewritten, reordered or removed.\par\medskip
\noindent\textit{Editorial note.} Following the multi-claim convention for this source, source\_url is the source-level URL https://stacks.math.columbia.edu/ and the source label lemma-closed-subscheme-proj-finite-type is carried in claim\_key and here rather than in a per-claim tag URL.\par\medskip
\noindent\textit{Editorial note.} The preamble is the chapter's shared Stacks Project preamble, retained as sources/stacks-divisors/retained-preamble.tex: a line-for-line excerpt of the shared Stacks preamble staged with this snapshot. No macro was invented or redefined, and this card adds no preamble text of its own.\par\medskip
\noindent\textit{Editorial note.} Statement span 8144-8162 and proof span 8164-8187 come from the halfA pool card operations/pool-cards/stacks-divisors-halfA.json and were re-read in the pinned chapter file before this card was written.\par\medskip
\noindent\textit{Editorial note.} The proof refers to the auxiliary local data introduced in the proof of the earlier companion statement and to an integer d fixed there; the author's notation is preserved and no letter was renamed.\par\medskip
\noindent\textit{Editorial note.} The proof contains no omissions or pointer-only steps; it is delivered as written.\par\medskip
\noindent\textit{Editorial note.} 3 cross-references occur inside the counted spans; they are listed in external\_references as navigation pointers only.\par\medskip
\noindent\textit{Allowed standard prerequisites (editorial).}\begin{itemize}
\item The description of an arbitrary closed subscheme of a relative Proj as a relative Proj of the graded quotient by an explicitly constructed kernel, cited from the same chapter and established by a previous card of this lane. Cited at proof lines 8164-8166, together with the auxiliary data and the integer d produced in the proof of its finite type companion. Supplies the graded ideal whose d-th piece has to be replaced by a finite type submodule with the same effect.
\item A quasi-coherent module over a quasi-compact quasi-separated scheme is the filtered union of its finite type quasi-coherent submodules, cited from the properties chapter. Cited at proof lines 8179-8181. Produces a single finite type submodule containing all the finitely many sections that have to be controlled.
\item A finite type module over a quasi-compact scheme is generated by finitely many sections over each member of a finite affine cover. Used at proof lines 8174-8178. Supplies the sections g over the members of a finite affine cover of the given open that have to lie in the submodule.
\end{itemize}\medskip
\ifdefined\hypertarget\hypertarget{sourcepointer-1}{}\fi
\noindent\textit{Source reference.} Original lemma 0801: focused-seeds/AGP/sources/stacks-snapshot-20261011/tags.txt, source lines 10166--10166. Same-chapter reference used in the retained author text at source line(s) 7991, 8111, 8166. It is supplied only as a navigation pointer so that the retained text compiles; the counted item delivers the author's own statement and the author's complete proof, and no part of the referenced item is copied in or substituted for the counted argument.\par\medskip
\ifdefined\hypertarget\hypertarget{sourcepointer-2}{}\fi
\noindent\textit{Source reference.} Original lemma 0802: focused-seeds/AGP/sources/stacks-snapshot-20261011/tags.txt, source lines 10167--10167. Same-chapter reference used in the retained author text at source line(s) 8140, 8169. It is supplied only as a navigation pointer so that the retained text compiles; the counted item delivers the author's own statement and the author's complete proof, and no part of the referenced item is copied in or substituted for the counted argument.\par\medskip
\ifdefined\hypertarget\hypertarget{sourcepointer-3}{}\fi
\noindent\textit{Source reference.} Original lemma 01PG: focused-seeds/AGP/sources/stacks-snapshot-20261011/tags.txt, source lines 2229--2229. Same-chapter reference used in the retained author text at source line(s) 8130, 8178. It is supplied only as a navigation pointer so that the retained text compiles; the counted item delivers the author's own statement and the author's complete proof, and no part of the referenced item is copied in or substituted for the counted argument.\par\medskip
\begin{lemma}
\label{lemma-closed-subscheme-proj-finite-type}
Let $S$ be a quasi-compact and quasi-separated scheme.
Let $\mathcal{A}$ be a quasi-coherent graded $\mathcal{O}_S$-algebra.
Let $p : X = \underline{\text{Proj}}_S(\mathcal{A}) \to S$ be the relative
Proj of $\mathcal{A}$. Let $i : Z \to X$ be a closed subscheme.
Let $U \subset S$ be an open. Assume that
\begin{enumerate}
\item $p$ is quasi-compact,
\item $i$ of finite presentation,
\item $U \cap p(i(Z)) = \emptyset$,
\item $U$ is quasi-compact,
\item $\mathcal{A}_n$ is a finite type $\mathcal{O}_S$-module for all $n$.
\end{enumerate}
Then there exists a $d > 0$ and a quasi-coherent finite type
$\mathcal{O}_S$-submodule $\mathcal{F} \subset \mathcal{A}_d$ with (a)
$Z = \underline{\text{Proj}}_S(\mathcal{A}/\mathcal{F}\mathcal{A})$
and (b) the support of $\mathcal{A}_d/\mathcal{F}$ is disjoint from $U$.
\end{lemma}

\begin{proof}
Let $\mathcal{I} \subset \mathcal{A}$ be the sheaf of quasi-coherent
graded ideals constructed in Lemma \ref{lemma-closed-subscheme-proj}.
Let $U_i$, $R_i$, $A_i$, $I_i$, $f_{i, j}$, $h_{i, j, k}$, and $d$
be as constructed in the proof of
Lemma \ref{lemma-closed-subscheme-proj-finite}.
Since $U \cap p(i(Z)) = \emptyset$ we see that
$\mathcal{I}_d|_U = \mathcal{A}_d|_U$ (by our construction of
$\mathcal{I}$ as a kernel). Since $U$ is quasi-compact we
can choose a finite affine open covering $U = W_1 \cup \ldots \cup W_m$.
Since $\mathcal{A}_d$ is of finite type we can find finitely many sections
$g_{t, s} \in \mathcal{A}_d(W_t)$ which generate
$\mathcal{A}_d|_{W_t} = \mathcal{I}_d|_{W_t}$
as an $\mathcal{O}_{W_t}$-module. To finish the proof, note that by
Properties, Lemma \ref{properties-lemma-quasi-coherent-colimit-finite-type}
there exists a finite type $\mathcal{F} \subset \mathcal{I}_d$
such that all the local sections
$$
h_{i, j, k}f_{i, j}^{(d - \deg(h_{i, j, k}))/\deg(f_{i, j})}
\quad\text{and}\quad
g_{t, s}
$$
are sections of $\mathcal{F}$. By construction $\mathcal{F}$ is a solution.
\end{proof}

\end{document}
